\scanpage{3}
\uline{Théorème} Soit $A$ un schéma abélien projectif sur $S$. Alors $\underline{\mathrm{Pic}}^{\tau}\deleted{\unsure{A|X}}_{A/S}\inserted{=\underline{\mathrm{Pic}}^{0}_{A/S}}$ est un schéma abélien sur $S$, ouvert et fermé dans $\underline{\mathrm{Pic}}_{A/S}$ ayant $=$ dimension relative sur $S$ que $A$.
% slanted
\margin{Si on ne veut pas se servir de la nullité de la torsion, on trouve que $\underline{\mathrm{Pic}}_{X/S}$ est simple sur $S$ en les pts de $\underline{\mathrm{Pic}}^{0}_{A/S}$, \deleted{\unsure{et on}} \inserted{et on} voit alors que $\varphi$ est un morphisme \uline{plat}, donc ouvert, donc $\underline{\mathrm{Pic}}^{0}_{A/S}$ est aussi ouvert dans $\underline{\mathrm{Pic}}_{A/S}$, c'est donc un schéma abélien sur $S$ par le stratagème \unsure{indiqué|ind.} \unsure{ci-dessous|ci-dessus} \ldots.}

Comme $\deleted{\unsure{X}}A$ est simple sur $S$ et que ses fibres \uline{ont des \unsure{Ner. Sév.|Nér. Sév.|Ner. Sev.} sans torsion}, on sait que $\underline{\mathrm{Pic}}^{\tau}_{\deleted{\unsure{X}}A/S}$ est propre sur $S$, et fermé dans $\underline{\mathrm{Pic}}_{A/S}$. Il est en tous cas \uline{ouvert} dans $\underline{\mathrm{Pic}}^{\tau}_{A/S}$.

Il \unsure{reste|veut} à montrer qu'il est \uline{simple} sur $S$ (on sait déjà que ses fibres sont universellement connexes), et a la bonne dimension.

Soit donc $\underline L$ un Module ample sur $A$ rel. $S$, il définit une section $s$ de $\underline{\mathrm{Pic}}_{A/S}$ sur $S$, et un \inserted{morphisme}
\deleted{\unsure{définit|désigne} par translation un sous-schéma ouvert et fermé}
\[
\deleted{P=s+\underline{\mathrm{Pic}}^{0}_{A/S}}
\]
\deleted{de $\underline{\mathrm{Pic}}_{A/S}$. Or on a un morphisme}
\[
\varphi_s:\ A\longrightarrow \deleted{\unsure{P}}A'\deleted{\unsure{0|o}}=\underline{\mathrm{Pic}}^\tau_{X/S}
\]
obtenu posant
\[
\varphi_s(t)=\text{classe de }\tau_{-t}^{*}(\underline L)\,\underline L^{-1}
\]
[Comme $\varphi_s$ transforme section unité en section unité, c'est un hom de schémas en groupes (\unsure{réf|rif})]. \deleted{\unsure{Donc}} L'hom $A_0\to A'_0$ défini par $\varphi$ est \uline{surjectif} [car le noyau est fini d'après le critère d'amplitude de Weil, d'autre part
\[
\dim A'_0\underset{\text{\unsure{déformation du fibre|définition du fibre|déformation du Pic|déformation de Pic}}}{\underbrace{\leq}} \dim H^1(A_0,\underline O_{A_0})
\underset{\substack{\text{th. de Hopf }+\\ \text{nullité de }H^i(A_0,\underline O_{A_0})\text{ pour }i>\dim A_0}}{\underbrace{\leq}}\dim A_0
\]
% archive sheet number: 2

